\section{Trace adjoints and rectangular Kraus maps}

\begin{definition}[Trace-pairing adjoint]\label{def:trace_pairing_adjoint}
    \lean{Matrix.traceAdjointMap}
    \leanok
    The trace-pairing adjoint $E^* : \MN{D'} \to \MN{D}$ of a linear map
    $E : \MN{D} \to \MN{D'}$ between matrix algebras of possibly different
    dimensions is the adjoint for the bilinear pairing
    $(A,B)\mapsto \tr(AB)$.
\end{definition}

\begin{definition}[Completely positive map in rectangular Kraus form]
    \label{def:is_kraus_cp}
    \lean{IsKrausCP}
    \leanok
    A linear map $\mathcal S$ between matrix algebras is \emph{completely
    positive in rectangular Kraus form} if there are finitely many operators
    $A_i:H\to K$ such that
    \begin{align}
        \mathcal S(X)&=\sum_i A_iXA_i^\dagger.
        \notag
    \end{align}
    No trace-preservation normalization is imposed.
\end{definition}

\begin{lemma}[Agreement with square-map complete positivity]
    \label{lem:is_kraus_cp_iff_is_cp_map}
    \lean{isKrausCP_iff_isCPMap}
    \leanok
    \uses{def:is_kraus_cp, def:cp_map}
    When the input and output matrix algebras agree, rectangular Kraus complete
    positivity is equivalent to the square-map notion of complete positivity.
\end{lemma}

\begin{lemma}[Basic properties of rectangular Kraus complete positivity]
    \label{lem:is_kraus_cp_basic}
    \lean{IsKrausCP.map_posSemidef, IsKrausCPTP.isKrausCP}
    \leanok
    \uses{def:is_kraus_cp, def:is_kraus_cptp}
    A map in rectangular Kraus form sends positive semidefinite matrices to
    positive semidefinite matrices. Every trace-preserving completely positive
    Kraus map is a completely positive Kraus map.
\end{lemma}

\begin{lemma}[Sum of rectangular Kraus maps]
    \label{lem:is_kraus_cp_add}
    \lean{IsKrausCP.add}
    \leanok
    \uses{def:is_kraus_cp}
    The sum of two completely positive maps in rectangular Kraus form is
    completely positive in rectangular Kraus form.
\end{lemma}

\begin{proof}\leanok
    Concatenating Kraus families for the two summands gives a Kraus family for
    their sum.
\end{proof}

\begin{definition}[Trace-preserving completely positive map]
    \label{def:is_kraus_cptp}
    \lean{IsKrausCPTP}
    \leanok
    A linear map $\mathcal{S}$ between matrix algebras is \emph{trace-preserving
        completely positive} if it has a Kraus form
    $\mathcal{S}(X)=\sum_i A_iXA_i^\dagger$ with
    $\sum_i A_i^\dagger A_i=I$. The Kraus operators may be rectangular, so the
    input and output dimensions need not agree.
\end{definition}

\begin{theorem}[Trace preservation completes a Kraus channel]
    \label{thm:is_kraus_cptp_of_is_kraus_cp_trace_preserving}
    \lean{isKrausCPTP_of_isKrausCP_trace_preserving}
    \leanok
    \uses{def:is_kraus_cp, def:is_kraus_cptp}
    A completely positive map in rectangular Kraus form that preserves the
    matrix trace is trace-preserving completely positive.
\end{theorem}

\begin{proof}\leanok
    If $(A_i)_i$ is a Kraus family, trace preservation and cyclicity give
    $\tr((\sum_i A_i^\dagger A_i)X)=\tr(X)$ for every $X$.
    Nondegeneracy of the trace pairing implies
    $\sum_i A_i^\dagger A_i=\Id$.
\end{proof}

\begin{theorem}[Trace preservation from a Kraus resolution]
    \label{thm:is_kraus_cptp_trace_map}
    \lean{IsKrausCPTP.trace_map}
    \leanok
    \uses{def:is_kraus_cptp}
    Every trace-preserving completely positive map in rectangular Kraus form
    preserves the matrix trace.
\end{theorem}

\begin{proof}
    \leanok
    If $\mathcal S(X)=\sum_a A_aXA_a^\dagger$ and
    $\sum_a A_a^\dagger A_a=I$, cyclicity gives
    \begin{align}
        \tr\mathcal S(X)
        &=\sum_a\tr(A_aXA_a^\dagger)
          =\tr\left(\left(\sum_a A_a^\dagger A_a\right)X\right)
          =\tr X.
        \notag
    \end{align}
\end{proof}

\begin{theorem}[The trace adjoint of a channel is unital]
    \label{thm:is_kraus_cptp_trace_adjoint_one}
    \lean{IsKrausCPTP.traceAdjointMap_one}
    \leanok
    \uses{def:is_kraus_cptp, def:trace_pairing_adjoint}
    If $\mathcal S:\MN{D}\to\MN{D'}$ is trace-preserving and completely
    positive, then its trace-pairing adjoint satisfies
    \begin{align}
        \mathcal S^*(\Id_{D'})
        &=\Id_D.
        \notag
    \end{align}
    This is the Schrödinger--Heisenberg duality of
    \cite[Section~1.2]{Wolf2012Quantum}.
\end{theorem}

\begin{proof}
    \leanok
    \uses{thm:trace_pairing_adjoint_identity,
        thm:is_kraus_cptp_trace_map, thm:trace_nondeg}
    For every $X\in\MN{D}$, trace duality and trace preservation give
    \begin{align}
        \tr\bigl(\mathcal S^*(\Id_{D'})X\bigr)
        &=\tr\bigl(\mathcal S(X)\bigr)
        =\tr(X)
        =\tr(\Id_DX).
        \notag
    \end{align}
    Nondegeneracy of the trace pairing proves the identity.
\end{proof}

\begin{lemma}[Positivity preservation for Kraus maps]
    \label{lem:is_kraus_cptp_map_pos}
    \lean{IsKrausCPTP.map_posSemidef}
    \leanok
    \uses{def:is_kraus_cptp}
    If $\mathcal S$ is trace-preserving completely positive and $X\geq0$, then
    $\mathcal S(X)\geq0$.
\end{lemma}

\section{The Choi--Jamiolkowski representation of maps between matrix algebras}

Fix $d,d'\geq 1$. For a linear map $T:\MN{d}\to\MN{d'}$ between matrix
algebras of possibly different dimensions, write $\Omega_d$ for the maximally
entangled vector on the \emph{input} factor and $\tau$ for the associated
operator on $\C^{d'}\otimes\C^{d}$, with the output factor first. The partial
trace over the output factor is written $\tr_A$ and the partial trace over
the input factor is written $\tr_B$.

\begin{definition}[Choi matrix of a map between matrix algebras]
    \label{def:choi_matrix_rect}
    \lean{ChoiRectangular.choiMatrix, ChoiRectangular.mapOfChoiMatrix}
    \leanok
    \uses{def:maximally_entangled, def:tensor_map_id}
    The \emph{Choi matrix} of a linear map $T:\MN{d}\to\MN{d'}$ is
    \begin{align}
        \tau &= (T\otimes\operatorname{id}_d)
          \bigl(\ketbra{\Omega_d}{\Omega_d}\bigr),
        \notag
    \end{align}
    an operator on $\C^{d'}\otimes\C^{d}$. This is the correspondence of
    \cite[Proposition~2.1]{Wolf2012Quantum}.
\end{definition}

\begin{theorem}[Mutual inverses of the Choi correspondence]
    \label{thm:choi_rect_mutual_inverse}
    \lean{ChoiRectangular.mapOfChoiMatrix_choiMatrix,
        ChoiRectangular.choiMatrix_mapOfChoiMatrix,
        ChoiRectangular.choiMatrix_injective}
    \leanok
    \uses{def:choi_matrix_rect}
    The assignments $T\mapsto\tau$ and $\tau\mapsto T$, the latter defined
    entrywise by
    \begin{align}
        T(B)_{i_1 j_1}
          &= d\sum_{i_2,j_2=0}^{d-1}
            \tau_{(i_1,i_2),(j_1,j_2)}\,B_{i_2 j_2},
        \notag
    \end{align}
    are mutual inverses: every linear map is recovered from its Choi matrix,
    and every operator on $\C^{d'}\otimes\C^{d}$ is the Choi matrix of the
    map it defines. This is \cite[Proposition~2.1, Equation~(2.4)]{Wolf2012Quantum}.
\end{theorem}

\begin{proof}
    \leanok
    The $(i_2,j_2)$-slice of $\ketbra{\Omega_d}{\Omega_d}$ is the matrix
    unit $E_{i_2 j_2}$ scaled by $1/d$; substituting it into the defining
    entries and using linearity gives both compositions.
\end{proof}

\begin{theorem}[Trace-pairing formula]
    \label{thm:choi_rect_trace_pairing}
    \lean{ChoiRectangular.trace_pairing}
    \leanok
    \uses{def:choi_matrix_rect, thm:choi_rect_mutual_inverse}
    For every $A\in\MN{d'}$ and $B\in\MN{d}$,
    \begin{align}
        \tr\!\bigl(A\,T(B)\bigr)
          &= d\,\tr\!\bigl(\tau\,(A\otimes B^{T})\bigr).
        \notag
    \end{align}
    This is the second line of \cite[Proposition~2.1, Equation~(2.1)]{Wolf2012Quantum}.
\end{theorem}

\begin{proof}
    \leanok
    Expand $\tr(A\,T(B))$ entrywise, insert the inverse assignment of
    Theorem~\ref{thm:choi_rect_mutual_inverse}, and identify the resulting
    four-fold sum with $d\,\tr(\tau\,(A\otimes B^{T}))$.
\end{proof}

\begin{theorem}[Hermiticity preservation]
    \label{thm:choi_rect_hermiticity}
    \lean{ChoiRectangular.choiMatrix_isHermitian_iff_hermiticityPreserving}
    \leanok
    \uses{def:choi_matrix_rect}
    The Choi matrix is Hermitian, $\tau=\tau^\dagger$, if and only if
    $T(B^\dagger)=T(B)^\dagger$ for every $B\in\MN{d}$. This is the
    Hermiticity clause of \cite[Proposition~2.1]{Wolf2012Quantum}.
\end{theorem}

\begin{proof}
    \leanok
    The entries of $\tau$ along the diagonal slices are the entries of
    $T(E_{i_2 j_2})$ scaled by $1/d$, so Hermiticity of $\tau$ is the
    relation $T(E_{j_2 i_2})=T(E_{i_2 j_2})^\dagger$ on matrix units; both
    directions then follow by linearity.
\end{proof}

\begin{theorem}[Complete positivity]
    \label{thm:choi_rect_cp}
    \lean{ChoiRectangular.isKrausCP_iff_choiMatrix_posSemidef,
        ChoiRectangular.exists_isKrausCP_of_posSemidef,
        ChoiRectangular.choiMatrix_of_kraus_posSemidef,
        ChoiRectangular.krausOfChoiDecomp, ChoiRectangular.krausOfChoiDecomp_spec,
        ChoiRectangular.exists_kraus_of_choiMatrix_eq_sum_vecMulVec}
    \leanok
    \uses{def:choi_matrix_rect, def:is_kraus_cp}
    The map $T$ is completely positive, in the sense that it admits a Kraus
    representation $T(X)=\sum_{j}K_jXK_j^\dagger$ with
    $K_j\in M_{d'\times d}(\C)$, if and only if $\tau\geq0$. This is the
    complete-positivity clause of \cite[Proposition~2.1]{Wolf2012Quantum}.
\end{theorem}

\begin{proof}
    \leanok
    If $T(X)=\sum_j K_jXK_j^\dagger$, then
    $\tau=\sum_j\ketbra{\psi_j}{\psi_j}$ with
    $(\psi_j)_{(i_1,i_2)}=(K_j)_{i_1 i_2}/\sqrt{d}$, so $\tau\geq0$.
    Conversely, a positive semidefinite $\tau$ decomposes into rank-one
    outer products $\sum_j\ketbra{\psi_j}{\psi_j}$, and rescaling the
    vectors $\psi_j$ by $\sqrt{d}$ gives Kraus operators via the inverse
    assignment.
\end{proof}

\begin{theorem}[Unitality]
    \label{thm:choi_rect_unital}
    \lean{ChoiRectangular.traceRight_choiMatrix,
        ChoiRectangular.unital_iff_traceRight_choiMatrix}
    \leanok
    \uses{def:choi_matrix_rect, def:partial_trace}
    The input-factor partial trace of the Choi matrix is
    $\tr_B(\tau)=T(\Id_d)/d$; hence $T(\Id_d)=\Id_{d'}$ if and only if
    $\tr_B(\tau)=\Id_{d'}/d$. This is the unitality clause of
    \cite[Proposition~2.1]{Wolf2012Quantum}.
\end{theorem}

\begin{proof}
    \leanok
    The diagonal slices of $\ketbra{\Omega_d}{\Omega_d}$ sum to $\Id_d/d$,
    so linearity gives $\tr_B(\tau)=T(\Id_d)/d$; cancelling the factor
    $1/d$ gives the equivalence.
\end{proof}

\begin{theorem}[Preservation of the trace]
    \label{thm:choi_rect_tp}
    \lean{ChoiRectangular.traceLeft_choiMatrix,
        ChoiRectangular.traceAdjointMap_one_eq_one_iff_tracePreserving,
        ChoiRectangular.tracePreserving_iff_traceLeft_choiMatrix}
    \leanok
    \uses{def:choi_matrix_rect, def:partial_trace, def:trace_pairing_adjoint}
    The output-factor partial trace of the Choi matrix is
    $\tr_A(\tau)=(T^*(\Id_{d'}))^{T}/d$, where $T^*$ is the trace-pairing
    adjoint; hence $T^*(\Id_{d'})=\Id_d$, equivalently $T$ preserves the
    trace, if and only if $\tr_A(\tau)=\Id_d/d$. This is the
    trace-preservation clause of \cite[Proposition~2.1]{Wolf2012Quantum}.
\end{theorem}

\begin{proof}
    \leanok
    The output-factor partial trace has entries
    $\tr(T(E_{i_2 j_2}))/d$, which are the entries of
    $(T^*(\Id_{d'}))^{T}/d$ by the defining trace pairing of the adjoint.
    The equivalence $\tr(T(X))=\tr(X)\Leftrightarrow T^*(\Id_{d'})=\Id_d$
    is the nondegeneracy of the trace pairing, and transposition and the
    factor $1/d$ cancel.
\end{proof}

\begin{theorem}[Trace normalization]
    \label{thm:choi_rect_normalization}
    \lean{ChoiRectangular.trace_choiMatrix}
    \leanok
    \uses{def:choi_matrix_rect, thm:choi_rect_tp}
    The full trace of the Choi matrix is
    $\tr(\tau)=\tr(T^*(\Id_{d'}))/d$. This is the normalization clause of
    \cite[Proposition~2.1]{Wolf2012Quantum}.
\end{theorem}

\begin{proof}
    \leanok
    Take the trace of the identity
    $\tr_A(\tau)=(T^*(\Id_{d'}))^{T}/d$ of Theorem~\ref{thm:choi_rect_tp}.
\end{proof}

\begin{theorem}[Doubly-stochastic maps]
    \label{thm:choi_rect_doubly_stochastic}
    \lean{ChoiRectangular.doublyStochastic_iff_partialTraces_proportional}
    \leanok
    \uses{def:choi_matrix_rect, thm:choi_rect_unital, thm:choi_rect_tp}
    The conditions $T(\Id_d)\propto\Id_{d'}$ and
    $T^*(\Id_{d'})\propto\Id_d$ hold if and only if
    $\tr_B(\tau)\propto\Id_{d'}$ and $\tr_A(\tau)\propto\Id_d$. This is the
    doubly-stochastic clause of \cite[Proposition~2.1]{Wolf2012Quantum}.
\end{theorem}

\begin{proof}
    \leanok
    Immediate from the partial-trace identities
    $\tr_B(\tau)=T(\Id_d)/d$ and $\tr_A(\tau)=(T^*(\Id_{d'}))^{T}/d$ of
    Theorems~\ref{thm:choi_rect_unital} and~\ref{thm:choi_rect_tp}.
\end{proof}

\section{The Kraus representation theorem for rectangular matrix algebras}

\begin{theorem}[Kraus normalization conditions]
    \label{thm:kraus_rect_normalization}
    \lean{ChoiRectangular.kraus_tp_iff_sum_conjTranspose_mul,
        ChoiRectangular.kraus_unital_iff_sum_mul_conjTranspose,
        ChoiRectangular.kraus_sum_conjTranspose_mul_of_tracePreserving}
    \leanok
    \uses{def:is_kraus_cp}
    Let $T(X)=\sum_{j=1}^{r}K_jXK_j^\dagger$ with
    $K_j\in M_{d'\times d}(\C)$. Then $T$ is trace preserving if and only
    if $\sum_j K_j^\dagger K_j=\Id_d$, and unital if and only if
    $\sum_j K_jK_j^\dagger=\Id_{d'}$. This is item~1 of
    \cite[Theorem~2.1]{Wolf2012Quantum}.
\end{theorem}

\begin{proof}
    \leanok
    Cyclicity of the trace gives
    $\tr\bigl(\bigl(\sum_j K_j^\dagger K_j\bigr)X\bigr)=\tr(T(X))$ for every
    $X$, so the trace-preserving equivalence follows from nondegeneracy of
    the trace pairing; the unital equivalence is the evaluation at $X=\Id_d$.
\end{proof}

\begin{definition}[Kraus cardinality and Kraus rank]
    \label{def:kraus_rank_rect}
    \lean{ChoiRectangular.HasKrausCard, ChoiRectangular.HasKrausRankLE,
        ChoiRectangular.choiRank, ChoiRectangular.hasKrausCard_mono}
    \leanok
    \uses{def:is_kraus_cp, def:choi_matrix_rect}
    A map $T:\MN{d}\to\MN{d'}$ has \emph{Kraus cardinality} $r$ if it has an
    exact $r$-operator Kraus representation. Its \emph{Kraus rank} (Choi
    rank) is the rank of its Choi matrix, $r=\operatorname{rank}(\tau)$;
    following \cite[Theorem~2.1, footnote]{Wolf2012Quantum}, this is
    distinguished from the rank of $T$ as a linear map.
\end{definition}

\begin{theorem}[Minimality of the Kraus rank]
    \label{thm:kraus_rect_rank_minimal}
    \lean{ChoiRectangular.hasKrausCard_choiRank_of_isKrausCP,
        ChoiRectangular.choiRank_le_of_hasKrausCard,
        ChoiRectangular.choiRank_isLeast_hasKrausCard_of_isKrausCP,
        ChoiRectangular.choiRank_le_mul}
    \leanok
    \uses{def:kraus_rank_rect, thm:choi_rect_cp}
    The minimal number of Kraus operators of a completely positive map
    $T:\MN{d}\to\MN{d'}$ is $r=\operatorname{rank}(\tau)\leq dd'$. This is
    item~2 of \cite[Theorem~2.1]{Wolf2012Quantum}.
\end{theorem}

\begin{proof}
    \leanok
    \uses{thm:choi_rect_cp}
    Every $r$-operator Kraus family gives
    $\tau=\sum_{j=1}^{r}\ketbra{\psi_j}{\psi_j}$, so
    $\operatorname{rank}(\tau)\leq r$. Conversely, the spectral
    decomposition of $\tau\geq0$ into $\operatorname{rank}(\tau)$ rank-one
    outer products yields a Kraus family of exactly
    $\operatorname{rank}(\tau)$ operators by
    Theorem~\ref{thm:choi_rect_cp}, and
    $\operatorname{rank}(\tau)\leq dd'$ since $\tau$ acts on a
    $dd'$-dimensional space.
\end{proof}

\begin{theorem}[Orthogonal minimal Kraus family]
    \label{thm:kraus_rect_orthogonal}
    \lean{ChoiRectangular.exists_kraus_orthogonal_of_isKrausCP}
    \leanok
    \uses{def:kraus_rank_rect, thm:kraus_rect_rank_minimal}
    Every completely positive map $T:\MN{d}\to\MN{d'}$ admits a Kraus
    representation with $r=\operatorname{rank}(\tau)$ Hilbert--Schmidt
    orthogonal Kraus operators, $\tr[K_i^\dagger K_j]\propto\delta_{ij}$;
    the off-diagonal traces vanish and the diagonal traces are nonzero.
    This is item~3 of \cite[Theorem~2.1]{Wolf2012Quantum}.
\end{theorem}

\begin{proof}
    \leanok
    \uses{thm:kraus_rect_rank_minimal}
    Choose the $\psi_j$ in the spectral decomposition of
    Theorem~\ref{thm:kraus_rect_rank_minimal} to be orthogonal
    eigenvectors; the corresponding Kraus operators then satisfy
    $\tr[K_i^\dagger K_j]=d\,\lambda_i\,\delta_{ij}$, where $\lambda_i$ is
    the $i$-th nonzero eigenvalue of $\tau$.
\end{proof}

\begin{theorem}[Unitary freedom of Kraus representations]
    \label{thm:kraus_rect_freedom}
    \lean{ChoiRectangular.kraus_isometry_freedom_iff,
        ChoiRectangular.kraus_unitary_freedom_iff,
        ChoiRectangular.kraus_isometry_freedom,
        ChoiRectangular.kraus_same_map_of_isometry_combination,
        ChoiRectangular.kraus_dual_eq_of_map_eq,
        ChoiRectangular.kraus_conjTranspose_mul_eq_of_map_eq}
    \leanok
    \uses{def:is_kraus_cp}
    Two sets of Kraus operators $\{K_j\}$ and $\{\widetilde{K}_\ell\}$
    represent the same completely positive map if and only if there is a
    unitary $U$ with $K_j=\sum_\ell U_{j\ell}\widetilde{K}_\ell$, where the
    smaller set is padded with zero operators; equivalently, after padding,
    the families are related by an isometry $V$ with $V^\dagger V=\Id$.
    This is item~4 of \cite[Theorem~2.1]{Wolf2012Quantum}.
\end{theorem}

\begin{proof}
    \leanok
    An isometric combination preserves the map by the orthogonality
    relations of the mixing matrix. Conversely, equal maps have equal
    trace-pairing adjoints, hence equal Gram matrices
    for the vectorized Kraus operators $(K_j)_{ab}$; equal Gram matrices
    are related by a unitary, whose rectangular submatrix is the claimed
    isometry after zero-padding.
\end{proof}
